% ============================================================================
% regulamin-inference-polska.tex -- Regulamin usług Inference Polska
%
% Rodzaj dokumentu: regulamin dla klientów (B2B), a nie praca badawcza w
% rozumieniu docs/research/00-RESEARCH-DOCUMENTATION.md. Nie wprowadza żadnych
% liczb ani faktów spoza tego pliku; jest jedynym źródłem treści regulaminu.
% Wspólna preambuła house-preamble.tex daje tylko wygląd.
%
% \paperstatus jest tu użyty na inne zastosowanie niż w szablonie badawczym:
% oznacza status wersji dokumentu (wersja robocza, do uzupełnienia danych).
%
% Dokument zawiera jeszcze pola w nawiasach [ ] do uzupełnienia przed
% publikacją. Build i publikacja PDF: .github/workflows/deploy-pages.yml.
% ============================================================================
\documentclass[11pt,a4paper]{article}
\newcommand{\paperstatus}{WERSJA ROBOCZA}
\newcommand{\headerleft}{Inference Polska --- Regulamin}
\newcommand{\pdftitleline}{Inference Polska: regulamin usług}
\newcommand{\pdfauthorline}{Konrad Cinkusz}
\def\houselang{polish}
% ============================================================================
% HOUSE-PREAMBLE.tex -- the estate's house LaTeX style, as a file you copy in
% and \input rather than a block you copy out.
%
% See 00-RESEARCH-DOCUMENTATION.md, "The house preamble is a file, not a
% block". Copy this file to <repo>/docs/papers/house-preamble.tex (or
% docs/research/papers/ for a study paper) and \input it from every document
% in the repository that wants the house look. One copy per repository; every
% document in that repository \inputs the same one.
%
% WHY A FILE. The estate's first attempt at this shipped a template whose
% preamble was meant to be copied into each new document. Copying is what
% drifts: `ab-ove`'s two overview editions carried none of the seven house
% markers -- not the colors, the section formatting, the header, the status
% marker, the path macro, the geometry or the link style -- and nothing said
% so, because a convention about copying correctly has nothing to fail. The
% Beamer theme, which was always distributed as a file, propagated
% byte-identically over the same period. Distribute the file.
%
% ----------------------------------------------------------------------------
% THE CONTRACT. The including document defines these BEFORE \input, and
% carries its own \documentclass and \begin{document} -- this file has
% neither, so that one preamble can serve documents with different classes.
%
%   \paperstatus    Right-hand header, and the title block's status line.
%                   A study paper: DRAFT or VERIFIED, mirroring the study's
%                   own header. Any other document: its own status, named in
%                   that document's header comment as repurposed.
%   \headerleft     Left-hand header. Usually "Project --- Document title".
%   \pdftitleline   PDF metadata title.
%   \pdfauthorline  PDF metadata author.
%
% Optional, both defaulted below if the document does not define them:
%
%   \houselang      babel language for this edition, e.g. \def\houselang{polish}.
%                   Defaults to english. Set it in the edition file, not here:
%                   it is the one thing that is genuinely per-edition.
%   \editionsuffix  This edition's diagram-file suffix, e.g.
%                   \def\editionsuffix{.pl}. Defaults to empty. \dgm reads it,
%                   so a figure is called by slug and resolves per edition --
%                   see "Figures" below.
%
% A minimal English document:
%
%   \documentclass[11pt,a4paper]{article}
%   \newcommand{\paperstatus}{DRAFT}
%   \newcommand{\headerleft}{Project --- Overview}
%   \newcommand{\pdftitleline}{Project: an overview}
%   \newcommand{\pdfauthorline}{Author Name}
%   % ============================================================================
% HOUSE-PREAMBLE.tex -- the estate's house LaTeX style, as a file you copy in
% and \input rather than a block you copy out.
%
% See 00-RESEARCH-DOCUMENTATION.md, "The house preamble is a file, not a
% block". Copy this file to <repo>/docs/papers/house-preamble.tex (or
% docs/research/papers/ for a study paper) and \input it from every document
% in the repository that wants the house look. One copy per repository; every
% document in that repository \inputs the same one.
%
% WHY A FILE. The estate's first attempt at this shipped a template whose
% preamble was meant to be copied into each new document. Copying is what
% drifts: `ab-ove`'s two overview editions carried none of the seven house
% markers -- not the colors, the section formatting, the header, the status
% marker, the path macro, the geometry or the link style -- and nothing said
% so, because a convention about copying correctly has nothing to fail. The
% Beamer theme, which was always distributed as a file, propagated
% byte-identically over the same period. Distribute the file.
%
% ----------------------------------------------------------------------------
% THE CONTRACT. The including document defines these BEFORE \input, and
% carries its own \documentclass and \begin{document} -- this file has
% neither, so that one preamble can serve documents with different classes.
%
%   \paperstatus    Right-hand header, and the title block's status line.
%                   A study paper: DRAFT or VERIFIED, mirroring the study's
%                   own header. Any other document: its own status, named in
%                   that document's header comment as repurposed.
%   \headerleft     Left-hand header. Usually "Project --- Document title".
%   \pdftitleline   PDF metadata title.
%   \pdfauthorline  PDF metadata author.
%
% Optional, both defaulted below if the document does not define them:
%
%   \houselang      babel language for this edition, e.g. \def\houselang{polish}.
%                   Defaults to english. Set it in the edition file, not here:
%                   it is the one thing that is genuinely per-edition.
%   \editionsuffix  This edition's diagram-file suffix, e.g.
%                   \def\editionsuffix{.pl}. Defaults to empty. \dgm reads it,
%                   so a figure is called by slug and resolves per edition --
%                   see "Figures" below.
%
% A minimal English document:
%
%   \documentclass[11pt,a4paper]{article}
%   \newcommand{\paperstatus}{DRAFT}
%   \newcommand{\headerleft}{Project --- Overview}
%   \newcommand{\pdftitleline}{Project: an overview}
%   \newcommand{\pdfauthorline}{Author Name}
%   % ============================================================================
% HOUSE-PREAMBLE.tex -- the estate's house LaTeX style, as a file you copy in
% and \input rather than a block you copy out.
%
% See 00-RESEARCH-DOCUMENTATION.md, "The house preamble is a file, not a
% block". Copy this file to <repo>/docs/papers/house-preamble.tex (or
% docs/research/papers/ for a study paper) and \input it from every document
% in the repository that wants the house look. One copy per repository; every
% document in that repository \inputs the same one.
%
% WHY A FILE. The estate's first attempt at this shipped a template whose
% preamble was meant to be copied into each new document. Copying is what
% drifts: `ab-ove`'s two overview editions carried none of the seven house
% markers -- not the colors, the section formatting, the header, the status
% marker, the path macro, the geometry or the link style -- and nothing said
% so, because a convention about copying correctly has nothing to fail. The
% Beamer theme, which was always distributed as a file, propagated
% byte-identically over the same period. Distribute the file.
%
% ----------------------------------------------------------------------------
% THE CONTRACT. The including document defines these BEFORE \input, and
% carries its own \documentclass and \begin{document} -- this file has
% neither, so that one preamble can serve documents with different classes.
%
%   \paperstatus    Right-hand header, and the title block's status line.
%                   A study paper: DRAFT or VERIFIED, mirroring the study's
%                   own header. Any other document: its own status, named in
%                   that document's header comment as repurposed.
%   \headerleft     Left-hand header. Usually "Project --- Document title".
%   \pdftitleline   PDF metadata title.
%   \pdfauthorline  PDF metadata author.
%
% Optional, both defaulted below if the document does not define them:
%
%   \houselang      babel language for this edition, e.g. \def\houselang{polish}.
%                   Defaults to english. Set it in the edition file, not here:
%                   it is the one thing that is genuinely per-edition.
%   \editionsuffix  This edition's diagram-file suffix, e.g.
%                   \def\editionsuffix{.pl}. Defaults to empty. \dgm reads it,
%                   so a figure is called by slug and resolves per edition --
%                   see "Figures" below.
%
% A minimal English document:
%
%   \documentclass[11pt,a4paper]{article}
%   \newcommand{\paperstatus}{DRAFT}
%   \newcommand{\headerleft}{Project --- Overview}
%   \newcommand{\pdftitleline}{Project: an overview}
%   \newcommand{\pdfauthorline}{Author Name}
%   \input{house-preamble}
%   \begin{document}
%   ...
%
% Its Polish sibling differs by two lines, both before the \input:
%
%   \def\houselang{polish}
%   \def\editionsuffix{.pl}
%
% ----------------------------------------------------------------------------
% PACKAGES THIS FILE DOES NOT LOAD. The house set below is what every document
% needs. A document that needs more -- tikz, pifont, longtable, fancyvrb,
% amssymb -- loads it in its own preamble AFTER the \input, and says in a
% comment that it is an addition rather than part of the house set, so the
% next reader can tell the two apart at a glance.
% ============================================================================

% ----------------------------------------------------------------------------
% Encoding, then language. fontenc before inputenc before babel: babel reads
% the active font encoding when it sets up its shorthands, so loading it first
% gives an edition that is translated but typeset wrong.
% ----------------------------------------------------------------------------
\usepackage[T1]{fontenc}
\usepackage[utf8]{inputenc}

% The edition's language, defaulted so a monolingual document defines nothing.
% \PassOptionsToPackage rather than \usepackage[\houselang]{babel}: the option
% list here is a macro, and PassOptionsToPackage is the mechanism built for
% passing one.
\providecommand{\houselang}{english}
\PassOptionsToPackage{\houselang}{babel}
\usepackage{babel}

% ----------------------------------------------------------------------------
% The house package set.
% ----------------------------------------------------------------------------
\usepackage{geometry}
\usepackage{xcolor}
\usepackage{hyperref}
\usepackage{graphicx}
\usepackage{enumitem}
\usepackage{fancyhdr}
\usepackage{titlesec}
\usepackage{parskip}
\usepackage{booktabs}
\usepackage{array}
\usepackage{tabularx}

% ----------------------------------------------------------------------------
% Page geometry.
% ----------------------------------------------------------------------------
\geometry{
    a4paper,
    left=2.2cm,
    right=2.2cm,
    top=2.5cm,
    bottom=2.5cm
}

% ----------------------------------------------------------------------------
% House colors. Extracted from the estate's first formal documents
% (business_analysis.tex, pitch-deck-demium.tex) and unchanged since: a
% document from any repository should look like it came from the same shop.
% ----------------------------------------------------------------------------
\definecolor{primaryblue}{RGB}{41,128,185}
\definecolor{darkblue}{RGB}{31,78,121}
\definecolor{lightgray}{RGB}{236,240,241}
\definecolor{darkgray}{RGB}{52,73,94}
\definecolor{accentgreen}{RGB}{39,174,96}
\definecolor{accentorange}{RGB}{230,126,34}
\definecolor{accentred}{RGB}{231,76,60}

% ----------------------------------------------------------------------------
% Section formatting.
% ----------------------------------------------------------------------------
\titleformat{\section}
    {\color{primaryblue}\Large\bfseries\sffamily}
    {\thesection.}{0.5em}{}[\titlerule]

\titleformat{\subsection}
    {\color{darkblue}\large\bfseries\sffamily}
    {\thesubsection}{0.5em}{}

\titleformat{\subsubsection}
    {\color{darkgray}\normalsize\bfseries\sffamily}
    {\thesubsubsection}{0.5em}{}

% ----------------------------------------------------------------------------
% Header and footer. \paperstatus rides in the right-hand header of every
% page, so a draft cannot be mistaken for a final document at any page a
% reader happens to open it at -- which is the whole reason the marker is in
% the header rather than only on the title page.
% ----------------------------------------------------------------------------
\pagestyle{fancy}
\fancyhf{}
\renewcommand{\headrulewidth}{0.4pt}
\fancyhead[L]{\small\sffamily\color{darkgray} \headerleft}
\fancyhead[R]{\small\sffamily\color{accentred} \paperstatus\ --- \today}
\fancyfoot[C]{\thepage}

% ----------------------------------------------------------------------------
% Links and PDF metadata.
% ----------------------------------------------------------------------------
\hypersetup{
    colorlinks=true,
    linkcolor=primaryblue,
    urlcolor=primaryblue,
    pdfauthor={\pdfauthorline},
    pdftitle={\pdftitleline}
}

% ----------------------------------------------------------------------------
% Figures. A diagram's source is one .mmd file rendered to vector PDF by the
% repository's render step; the document includes the rendered file.
%
% \dgm resolves the path from a slug and the edition, so the body text names a
% diagram the way the diagrams directory does and never writes a path:
%
%   \includegraphics[width=0.78\linewidth]{\dgm{a1-system-context}}
%
% resolves to ../diagrams/rendered/a1-system-context.pdf in the default
% edition and ../diagrams/rendered/a1-system-context.pl.pdf in an edition that
% set \def\editionsuffix{.pl}. A repository whose editions share one set of
% diagrams simply never sets \editionsuffix, and every edition resolves to the
% same file.
%
% Hardcoding the path instead is what the estate did first, and it costs one
% edit per figure per edition: `ab-ove` maintained `b1-reader-loop.pdf` in the
% English file against `b1-reader-loop.pl.pdf` in the Polish one, by hand, in
% both. At three figures that is survivable; `agent-eval-bench` has 23.
%
% \housediagramdir is overridable for a document that does not sit one level
% below the diagrams directory -- define it before the \input.
% ----------------------------------------------------------------------------
\providecommand{\editionsuffix}{}
\providecommand{\housediagramdir}{../diagrams/rendered}
\newcommand{\dgm}[1]{\housediagramdir/#1\editionsuffix.pdf}

% ----------------------------------------------------------------------------
% House convenience commands.
%
% \repopath is for a path, a filename or an identifier copied from the code;
% it is deliberately the only one of these, so that a document does not grow a
% private synonym for it. \finding marks the sentence a section exists to
% deliver.
% ----------------------------------------------------------------------------
\newcommand{\repopath}[1]{\colorbox{lightgray}{\texttt{#1}}}
\newcommand{\finding}[1]{{\color{primaryblue}\textbf{#1}}}

% ----------------------------------------------------------------------------
% The house title block, as a command so that every document's first page
% looks the same and none of them reimplements it.
%
%   \housetitle{Title}{Subtitle or empty}{Author}{Affiliation or repo URL}{Provenance line}
%
% The provenance line is the one this estate keeps paying for the absence of:
% a document that presents another document names it here, so a reader who
% finds the two disagreeing knows which one is the source. For a study paper
% it names the companion study and the commit; for a project overview, the
% markdown it presents.
% ----------------------------------------------------------------------------
\newcommand{\housetitle}[5]{%
    \thispagestyle{plain}%
    \begin{center}
        {\LARGE\bfseries\sffamily\color{darkblue} #1}

        \ifx\relax#2\relax\else
            \vspace{0.3cm}

            {\Large\sffamily\color{darkgray} #2}
        \fi

        \vspace{0.6cm}

        {\large #3}

        \vspace{0.2cm}

        {\normalsize\color{darkgray} #4}

        \vspace{0.2cm}

        {\normalsize\color{accentred}\sffamily \paperstatus\ --- \today\ --- #5}
    \end{center}

    \vspace{0.4cm}
}

%   \begin{document}
%   ...
%
% Its Polish sibling differs by two lines, both before the \input:
%
%   \def\houselang{polish}
%   \def\editionsuffix{.pl}
%
% ----------------------------------------------------------------------------
% PACKAGES THIS FILE DOES NOT LOAD. The house set below is what every document
% needs. A document that needs more -- tikz, pifont, longtable, fancyvrb,
% amssymb -- loads it in its own preamble AFTER the \input, and says in a
% comment that it is an addition rather than part of the house set, so the
% next reader can tell the two apart at a glance.
% ============================================================================

% ----------------------------------------------------------------------------
% Encoding, then language. fontenc before inputenc before babel: babel reads
% the active font encoding when it sets up its shorthands, so loading it first
% gives an edition that is translated but typeset wrong.
% ----------------------------------------------------------------------------
\usepackage[T1]{fontenc}
\usepackage[utf8]{inputenc}

% The edition's language, defaulted so a monolingual document defines nothing.
% \PassOptionsToPackage rather than \usepackage[\houselang]{babel}: the option
% list here is a macro, and PassOptionsToPackage is the mechanism built for
% passing one.
\providecommand{\houselang}{english}
\PassOptionsToPackage{\houselang}{babel}
\usepackage{babel}

% ----------------------------------------------------------------------------
% The house package set.
% ----------------------------------------------------------------------------
\usepackage{geometry}
\usepackage{xcolor}
\usepackage{hyperref}
\usepackage{graphicx}
\usepackage{enumitem}
\usepackage{fancyhdr}
\usepackage{titlesec}
\usepackage{parskip}
\usepackage{booktabs}
\usepackage{array}
\usepackage{tabularx}

% ----------------------------------------------------------------------------
% Page geometry.
% ----------------------------------------------------------------------------
\geometry{
    a4paper,
    left=2.2cm,
    right=2.2cm,
    top=2.5cm,
    bottom=2.5cm
}

% ----------------------------------------------------------------------------
% House colors. Extracted from the estate's first formal documents
% (business_analysis.tex, pitch-deck-demium.tex) and unchanged since: a
% document from any repository should look like it came from the same shop.
% ----------------------------------------------------------------------------
\definecolor{primaryblue}{RGB}{41,128,185}
\definecolor{darkblue}{RGB}{31,78,121}
\definecolor{lightgray}{RGB}{236,240,241}
\definecolor{darkgray}{RGB}{52,73,94}
\definecolor{accentgreen}{RGB}{39,174,96}
\definecolor{accentorange}{RGB}{230,126,34}
\definecolor{accentred}{RGB}{231,76,60}

% ----------------------------------------------------------------------------
% Section formatting.
% ----------------------------------------------------------------------------
\titleformat{\section}
    {\color{primaryblue}\Large\bfseries\sffamily}
    {\thesection.}{0.5em}{}[\titlerule]

\titleformat{\subsection}
    {\color{darkblue}\large\bfseries\sffamily}
    {\thesubsection}{0.5em}{}

\titleformat{\subsubsection}
    {\color{darkgray}\normalsize\bfseries\sffamily}
    {\thesubsubsection}{0.5em}{}

% ----------------------------------------------------------------------------
% Header and footer. \paperstatus rides in the right-hand header of every
% page, so a draft cannot be mistaken for a final document at any page a
% reader happens to open it at -- which is the whole reason the marker is in
% the header rather than only on the title page.
% ----------------------------------------------------------------------------
\pagestyle{fancy}
\fancyhf{}
\renewcommand{\headrulewidth}{0.4pt}
\fancyhead[L]{\small\sffamily\color{darkgray} \headerleft}
\fancyhead[R]{\small\sffamily\color{accentred} \paperstatus\ --- \today}
\fancyfoot[C]{\thepage}

% ----------------------------------------------------------------------------
% Links and PDF metadata.
% ----------------------------------------------------------------------------
\hypersetup{
    colorlinks=true,
    linkcolor=primaryblue,
    urlcolor=primaryblue,
    pdfauthor={\pdfauthorline},
    pdftitle={\pdftitleline}
}

% ----------------------------------------------------------------------------
% Figures. A diagram's source is one .mmd file rendered to vector PDF by the
% repository's render step; the document includes the rendered file.
%
% \dgm resolves the path from a slug and the edition, so the body text names a
% diagram the way the diagrams directory does and never writes a path:
%
%   \includegraphics[width=0.78\linewidth]{\dgm{a1-system-context}}
%
% resolves to ../diagrams/rendered/a1-system-context.pdf in the default
% edition and ../diagrams/rendered/a1-system-context.pl.pdf in an edition that
% set \def\editionsuffix{.pl}. A repository whose editions share one set of
% diagrams simply never sets \editionsuffix, and every edition resolves to the
% same file.
%
% Hardcoding the path instead is what the estate did first, and it costs one
% edit per figure per edition: `ab-ove` maintained `b1-reader-loop.pdf` in the
% English file against `b1-reader-loop.pl.pdf` in the Polish one, by hand, in
% both. At three figures that is survivable; `agent-eval-bench` has 23.
%
% \housediagramdir is overridable for a document that does not sit one level
% below the diagrams directory -- define it before the \input.
% ----------------------------------------------------------------------------
\providecommand{\editionsuffix}{}
\providecommand{\housediagramdir}{../diagrams/rendered}
\newcommand{\dgm}[1]{\housediagramdir/#1\editionsuffix.pdf}

% ----------------------------------------------------------------------------
% House convenience commands.
%
% \repopath is for a path, a filename or an identifier copied from the code;
% it is deliberately the only one of these, so that a document does not grow a
% private synonym for it. \finding marks the sentence a section exists to
% deliver.
% ----------------------------------------------------------------------------
\newcommand{\repopath}[1]{\colorbox{lightgray}{\texttt{#1}}}
\newcommand{\finding}[1]{{\color{primaryblue}\textbf{#1}}}

% ----------------------------------------------------------------------------
% The house title block, as a command so that every document's first page
% looks the same and none of them reimplements it.
%
%   \housetitle{Title}{Subtitle or empty}{Author}{Affiliation or repo URL}{Provenance line}
%
% The provenance line is the one this estate keeps paying for the absence of:
% a document that presents another document names it here, so a reader who
% finds the two disagreeing knows which one is the source. For a study paper
% it names the companion study and the commit; for a project overview, the
% markdown it presents.
% ----------------------------------------------------------------------------
\newcommand{\housetitle}[5]{%
    \thispagestyle{plain}%
    \begin{center}
        {\LARGE\bfseries\sffamily\color{darkblue} #1}

        \ifx\relax#2\relax\else
            \vspace{0.3cm}

            {\Large\sffamily\color{darkgray} #2}
        \fi

        \vspace{0.6cm}

        {\large #3}

        \vspace{0.2cm}

        {\normalsize\color{darkgray} #4}

        \vspace{0.2cm}

        {\normalsize\color{accentred}\sffamily \paperstatus\ --- \today\ --- #5}
    \end{center}

    \vspace{0.4cm}
}

%   \begin{document}
%   ...
%
% Its Polish sibling differs by two lines, both before the \input:
%
%   \def\houselang{polish}
%   \def\editionsuffix{.pl}
%
% ----------------------------------------------------------------------------
% PACKAGES THIS FILE DOES NOT LOAD. The house set below is what every document
% needs. A document that needs more -- tikz, pifont, longtable, fancyvrb,
% amssymb -- loads it in its own preamble AFTER the \input, and says in a
% comment that it is an addition rather than part of the house set, so the
% next reader can tell the two apart at a glance.
% ============================================================================

% ----------------------------------------------------------------------------
% Encoding, then language. fontenc before inputenc before babel: babel reads
% the active font encoding when it sets up its shorthands, so loading it first
% gives an edition that is translated but typeset wrong.
% ----------------------------------------------------------------------------
\usepackage[T1]{fontenc}
\usepackage[utf8]{inputenc}

% The edition's language, defaulted so a monolingual document defines nothing.
% \PassOptionsToPackage rather than \usepackage[\houselang]{babel}: the option
% list here is a macro, and PassOptionsToPackage is the mechanism built for
% passing one.
\providecommand{\houselang}{english}
\PassOptionsToPackage{\houselang}{babel}
\usepackage{babel}

% ----------------------------------------------------------------------------
% The house package set.
% ----------------------------------------------------------------------------
\usepackage{geometry}
\usepackage{xcolor}
\usepackage{hyperref}
\usepackage{graphicx}
\usepackage{enumitem}
\usepackage{fancyhdr}
\usepackage{titlesec}
\usepackage{parskip}
\usepackage{booktabs}
\usepackage{array}
\usepackage{tabularx}

% ----------------------------------------------------------------------------
% Page geometry.
% ----------------------------------------------------------------------------
\geometry{
    a4paper,
    left=2.2cm,
    right=2.2cm,
    top=2.5cm,
    bottom=2.5cm
}

% ----------------------------------------------------------------------------
% House colors. Extracted from the estate's first formal documents
% (business_analysis.tex, pitch-deck-demium.tex) and unchanged since: a
% document from any repository should look like it came from the same shop.
% ----------------------------------------------------------------------------
\definecolor{primaryblue}{RGB}{41,128,185}
\definecolor{darkblue}{RGB}{31,78,121}
\definecolor{lightgray}{RGB}{236,240,241}
\definecolor{darkgray}{RGB}{52,73,94}
\definecolor{accentgreen}{RGB}{39,174,96}
\definecolor{accentorange}{RGB}{230,126,34}
\definecolor{accentred}{RGB}{231,76,60}

% ----------------------------------------------------------------------------
% Section formatting.
% ----------------------------------------------------------------------------
\titleformat{\section}
    {\color{primaryblue}\Large\bfseries\sffamily}
    {\thesection.}{0.5em}{}[\titlerule]

\titleformat{\subsection}
    {\color{darkblue}\large\bfseries\sffamily}
    {\thesubsection}{0.5em}{}

\titleformat{\subsubsection}
    {\color{darkgray}\normalsize\bfseries\sffamily}
    {\thesubsubsection}{0.5em}{}

% ----------------------------------------------------------------------------
% Header and footer. \paperstatus rides in the right-hand header of every
% page, so a draft cannot be mistaken for a final document at any page a
% reader happens to open it at -- which is the whole reason the marker is in
% the header rather than only on the title page.
% ----------------------------------------------------------------------------
\pagestyle{fancy}
\fancyhf{}
\renewcommand{\headrulewidth}{0.4pt}
\fancyhead[L]{\small\sffamily\color{darkgray} \headerleft}
\fancyhead[R]{\small\sffamily\color{accentred} \paperstatus\ --- \today}
\fancyfoot[C]{\thepage}

% ----------------------------------------------------------------------------
% Links and PDF metadata.
% ----------------------------------------------------------------------------
\hypersetup{
    colorlinks=true,
    linkcolor=primaryblue,
    urlcolor=primaryblue,
    pdfauthor={\pdfauthorline},
    pdftitle={\pdftitleline}
}

% ----------------------------------------------------------------------------
% Figures. A diagram's source is one .mmd file rendered to vector PDF by the
% repository's render step; the document includes the rendered file.
%
% \dgm resolves the path from a slug and the edition, so the body text names a
% diagram the way the diagrams directory does and never writes a path:
%
%   \includegraphics[width=0.78\linewidth]{\dgm{a1-system-context}}
%
% resolves to ../diagrams/rendered/a1-system-context.pdf in the default
% edition and ../diagrams/rendered/a1-system-context.pl.pdf in an edition that
% set \def\editionsuffix{.pl}. A repository whose editions share one set of
% diagrams simply never sets \editionsuffix, and every edition resolves to the
% same file.
%
% Hardcoding the path instead is what the estate did first, and it costs one
% edit per figure per edition: `ab-ove` maintained `b1-reader-loop.pdf` in the
% English file against `b1-reader-loop.pl.pdf` in the Polish one, by hand, in
% both. At three figures that is survivable; `agent-eval-bench` has 23.
%
% \housediagramdir is overridable for a document that does not sit one level
% below the diagrams directory -- define it before the \input.
% ----------------------------------------------------------------------------
\providecommand{\editionsuffix}{}
\providecommand{\housediagramdir}{../diagrams/rendered}
\newcommand{\dgm}[1]{\housediagramdir/#1\editionsuffix.pdf}

% ----------------------------------------------------------------------------
% House convenience commands.
%
% \repopath is for a path, a filename or an identifier copied from the code;
% it is deliberately the only one of these, so that a document does not grow a
% private synonym for it. \finding marks the sentence a section exists to
% deliver.
% ----------------------------------------------------------------------------
\newcommand{\repopath}[1]{\colorbox{lightgray}{\texttt{#1}}}
\newcommand{\finding}[1]{{\color{primaryblue}\textbf{#1}}}

% ----------------------------------------------------------------------------
% The house title block, as a command so that every document's first page
% looks the same and none of them reimplements it.
%
%   \housetitle{Title}{Subtitle or empty}{Author}{Affiliation or repo URL}{Provenance line}
%
% The provenance line is the one this estate keeps paying for the absence of:
% a document that presents another document names it here, so a reader who
% finds the two disagreeing knows which one is the source. For a study paper
% it names the companion study and the commit; for a project overview, the
% markdown it presents.
% ----------------------------------------------------------------------------
\newcommand{\housetitle}[5]{%
    \thispagestyle{plain}%
    \begin{center}
        {\LARGE\bfseries\sffamily\color{darkblue} #1}

        \ifx\relax#2\relax\else
            \vspace{0.3cm}

            {\Large\sffamily\color{darkgray} #2}
        \fi

        \vspace{0.6cm}

        {\large #3}

        \vspace{0.2cm}

        {\normalsize\color{darkgray} #4}

        \vspace{0.2cm}

        {\normalsize\color{accentred}\sffamily \paperstatus\ --- \today\ --- #5}
    \end{center}

    \vspace{0.4cm}
}


\setlength{\parindent}{0pt}
\setlist[enumerate,1]{label=\arabic*.,leftmargin=*,itemsep=2pt}

\begin{document}

\housetitle{Regulamin usług Inference Polska}
  {Analiza spraw, rozmowy wstępne i kontakt, w trybie offline lub online}
  {Konrad Cinkusz, prowadzący działalność pod firmą Inference Polska}
  {[adres], NIP [numer], e-mail [adres]}
  {regulamin ogólny i zasady przetwarzania danych}

\section*{Część I. Regulamin ogólny usług i kontaktu}

\subsection*{§1 Postanowienia ogólne}
\begin{enumerate}
\item Regulamin określa zasady świadczenia usług przez Konrada Cinkusza, prowadzącego jednoosobową działalność gospodarczą pod nazwą Inference Polska, [adres], NIP [numer], e-mail [adres] (dalej: Usługodawca).
\item Usługi są kierowane wyłącznie do przedsiębiorców i osób wykonujących wolne zawody (np. gabinety, kancelarie, firmy) w związku z ich działalnością (B2B). Usługodawca nie zawiera umów z konsumentami.
\item Użyte pojęcia: Klient to podmiot zlecający usługę; Materiały to dokumenty, akta i dane przekazane przez Klienta do analizy; Analiza to wynik pracy Usługodawcy z użyciem modeli sztucznej inteligencji (AI); Tryb Fully private i Tryb Cloud opisano w Części IV i V; Rozmowa to spotkanie online w Google Meet.
\item Regulamin jest dostępny na stronie [adres strony] i wiąże Klienta po jego akceptacji w ofercie, e-mailu lub umowie. Indywidualna umowa lub oferta ma pierwszeństwo przed Regulaminem.
\end{enumerate}

\subsection*{§2 Zakres usług}
\begin{enumerate}
\item Usługodawca świadczy w szczególności: (a) analizę spraw z użyciem AI w trybie Fully private lub Cloud (Część IV i V), (b) wynajem urządzenia Private AI Node wraz z oprogramowaniem, na podstawie odrębnej umowy, (c) konsultacje i wdrożenia lokalnej AI.
\item Szczegółowy zakres, termin i cena każdego zlecenia wynikają z oferty zaakceptowanej przez Klienta.
\item Usługodawca może odmówić przyjęcia zlecenia, w szczególności gdy wymagałoby ono przetwarzania danych w sposób niezgodny z prawem lub z wybranym trybem.
\end{enumerate}

\subsection*{§3 Kontakt z Usługodawcą}
\begin{enumerate}
\item Kanały kontaktu: e-mail [adres], formularz na stronie [adres], telefon [numer], rozmowa w Google Meet po wcześniejszym umówieniu terminu.
\item Usługodawca odpowiada zwykle w ciągu [2] dni roboczych. Termin ma charakter orientacyjny i nie jest zobowiązaniem.
\item Kontakt, rozmowa wstępna i wycena nie tworzą umowy ani nie zobowiązują Klienta do zakupu.
\item Nie przesyłaj danych osobowych, danych o zdrowiu, tajemnicy zawodowej ani akt spraw przez formularz, zwykły e-mail ani komunikator przed uzgodnieniem trybu i sposobu ich przekazania (§ 9). W pierwszym kontakcie opisz sprawę ogólnie, bez danych identyfikujących osoby.
\item Usługodawca nie udziela pomocy w sytuacjach pilnych ani zagrażających życiu i zdrowiu; usługi nie zastępują służb ani lekarza.
\end{enumerate}

\subsection*{§4 Zawarcie umowy}
\begin{enumerate}
\item Proces: zapytanie, bezpłatna rozmowa wstępna, oferta przesłana e-mailem, jej akceptacja przez Klienta (e-mail wystarczy), a przy trybie Fully private i Cloud także podpisanie umowy powierzenia przetwarzania danych (§ 10), jeśli dane osobowe będą przetwarzane.
\item Oferta jest ważna [14] dni, o ile nie wskazano inaczej.
\item Usługodawca rozpoczyna pracę po akceptacji oferty, zawarciu umowy powierzenia (gdy jest potrzebna) i otrzymaniu płatności lub zaliczki, jeśli oferta ją przewiduje.
\end{enumerate}

\subsection*{§5 Rozmowa bezpłatna}
\begin{enumerate}
\item Rozmowa wstępna jest bezpłatna i służy ustaleniu potrzeb Klienta. Jej odbycie nie zobowiązuje żadnej ze Stron.
\item Rozmowy odbywają się w Google Meet. Usługodawca nagrywa je na zasadach z Części II. Klient może odmówić nagrywania, co nie wpływa na prawo do rozmowy.
\end{enumerate}

\section*{Część II. Rozmowy w Google Meet i nagrywanie}

\subsection*{§6 Nagrywanie rozmów}
\begin{enumerate}
\item Bezpłatne rozmowy z Usługodawcą (wstępne, rozpoznanie potrzeb, demonstracje) odbywają się w Google Meet i są nagrywane (obraz i dźwięk), aby Usługodawca mógł później sporządzić notatkę i wyciągnąć wnioski z pomocą AI.
\item Przed rozpoczęciem nagrania Usługodawca informuje o nim uczestników ustnie; Google Meet pokazuje też wskaźnik nagrywania. Informacja o nagrywaniu jest też w zaproszeniu na spotkanie.
\item Podstawą nagrywania jest zgoda uczestnika (art. 6 ust. 1 lit. a RODO), wyrażona przez pozostanie na rozmowie po poinformowaniu albo wprost na prośbę Usługodawcy; zgoda zapisuje się w nagraniu. Uczestnik może ją cofnąć w każdej chwili, co nie wpływa na zgodność z prawem wcześniejszego przetwarzania. Cofnięcie zgody oznacza usunięcie nagrania.
\item Uczestnik może odmówić nagrywania przed rozmową lub w jej trakcie. Wtedy rozmowa odbywa się bez nagrania, a Usługodawca robi tylko odręczne notatki. Odmowa nie wpływa na ofertę ani cenę.
\item Płatne sesje, analizy i spotkania robocze nie są nagrywane, chyba że Strony zgodzą się na to z góry i pisemnie.
\item Uczestnik występujący w imieniu Klienta informuje własnych pracowników lub współpracowników, którzy dołączają do rozmowy, o nagrywaniu.
\item Podczas bezpłatnych rozmów nie podawaj danych pacjentów, klientów ani treści akt. Przykłady opisuj w sposób anonimowy. Jeśli dane takie padły, uczestnik może zażądać usunięcia nagrania.
\end{enumerate}

\subsection*{§7 Przetwarzanie nagrań z użyciem AI}
\begin{enumerate}
\item Nagranie może zostać zamienione na transkrypcję i przeanalizowane przez modele AI wyłącznie w celu: sporządzenia notatki ze spotkania, przygotowania oferty, poprawy opisu usług i dokumentowania ustaleń.
\item [Do uzupełnienia: transkrypcja i analiza odbywają się lokalnie na sprzęcie Usługodawcy / z użyciem narzędzia X (odbiorca danych), przy czym dane mogą być przekazywane poza EOG na podstawie [mechanizm].] Aktualna lista odbiorców jest w Załączniku nr 1.
\item Nagrań i transkrypcji nie używa się do trenowania modeli ani nie przekazuje się ich osobom trzecim poza wskazanymi odbiorcami.
\item Nagranie jest przechowywane do [90] dni od rozmowy, po czym jest usuwane. Notatka bez danych identyfikujących osoby trzecie może być przechowywana dla celów ofertowych i rozliczeń do [12] miesięcy od ostatniego kontaktu, a w razie zawarcia umowy przez okres jej obowiązywania i przedawnienia roszczeń.
\item Wyniki pracy AI nie podejmują wobec uczestnika żadnych zautomatyzowanych decyzji w rozumieniu art. 22 RODO.
\item Uczestnikowi przysługuje prawo dostępu, sprostowania, usunięcia, ograniczenia, sprzeciwu, przeniesienia danych i skargi do Prezesa UODO. Kontakt: [e-mail]. Pełna klauzula jest w Załączniku nr 1.
\end{enumerate}

\section*{Część III. Usługa Analiza spraw: postanowienia wspólne}

\subsection*{§8 Opis i granice usługi}
\begin{enumerate}
\item Analiza spraw polega na przeanalizowaniu przez Usługodawcę, z użyciem modeli AI, Materiałów przekazanych przez Klienta (np. dokumentacji spraw, akt, pism, notatek) i przygotowaniu raportu. Raport może zawierać m.in. streszczenie, chronologię zdarzeń, wykaz kluczowych faktów i kwestii spornych, wykryte braki i niespójności oraz tabele zestawień. Zakres ustala oferta.
\item Usługa jest narzędziem wspomagającym pracę Klienta. Nie jest poradą prawną, medyczną, podatkową ani innym doradztwem i nie zastępuje oceny uprawnionego specjalisty po stronie Klienta. Usługa nie jest wyrobem medycznym i nie diagnozuje.
\item Modele AI mogą popełniać błędy, pomijać informacje lub podawać treści nieprawdziwe. Klient jest zobowiązany zweryfikować raport z źródłami (raport wskazuje, gdzie to możliwe, miejsce w Materiałach) przed jego użyciem i ponosi odpowiedzialność za decyzje podjęte na jego podstawie. Usługodawca nie gwarantuje określonego wyniku ani kompletności analizy.
\end{enumerate}

\subsection*{§9 Wybór trybu i przekazanie Materiałów}
\begin{enumerate}
\item Klient wybiera przed rozpoczęciem pracy jeden z trybów: Fully private (Część IV) albo Cloud (Część V). Wybór wpisuje się do oferty. Zmiana trybu w trakcie wymaga zgody obu Stron i może zmienić cenę i termin.
\item Materiały przekazuje się wyłącznie sposobem uzgodnionym w ofercie dla danego trybu. Materiałów nie wolno wysyłać zwykłym e-mailem ani przez komunikator.
\item Klient przed przekazaniem Materiałów usuwa z nich dane, które nie są potrzebne do analizy, i oznacza materiały zawierające dane szczególnych kategorii (np. o zdrowiu) lub objęte tajemnicą zawodową. Preferowane jest wcześniejsze zanonimizowanie lub pseudonimizacja.
\end{enumerate}

\subsection*{§10 Powierzenie przetwarzania danych}
\begin{enumerate}
\item Jeśli Materiały zawierają dane osobowe, Klient (administrator) powierza je Usługodawcy (podmiotowi przetwarzającemu) na podstawie umowy powierzenia z art. 28 RODO, zawartej przed przekazaniem Materiałów. Wzór jest w Załączniku nr 2 do wzoru umowy najmu Private AI Node i może być użyty odpowiednio.
\item Usługodawca przetwarza dane wyłącznie na udokumentowane polecenie Klienta, w zakresie niezbędnym do wykonania analizy.
\item Klient odpowiada za to, że ma podstawę prawną do przekazania Materiałów, w tym wobec osób, których dane dotyczą, oraz że zachowuje obowiązek tajemnicy zawodowej, jeśli go dotyczy, a przekazanie nie narusza przepisów o tej tajemnicy.
\end{enumerate}

\subsection*{§11 Obowiązki Klienta}
\begin{enumerate}
\item Klient przekazuje Materiały w czytelnej, uzgodnionej formie i opisuje cel analizy oraz pytania, na które ma odpowiedzieć raport.
\item Klient wskazuje osobę kontaktową i odpowiada na pytania Usługodawcy w terminie [3] dni roboczych. Zwłoka Klienta przesuwa termin wykonania.
\item Klient oświadcza, że ma prawo korzystać z Materiałów i że ich przekazanie nie narusza praw osób trzecich.
\end{enumerate}

\subsection*{§12 Raport, prawa i poufność}
\begin{enumerate}
\item Raport jest przekazywany w uzgodnionym formacie i kanale. Po zapłacie wynagrodzenia Klient może używać raportu na własny użytek zawodowy bez ograniczeń czasowych.
\item Usługodawca zachowuje w poufności Materiały i raport, również po zakończeniu umowy, bez ograniczeń w czasie. Nie używa ich do trenowania modeli ani w studiach przypadków bez odrębnej pisemnej zgody Klienta.
\item Klient zachowuje wszelkie prawa do swoich Materiałów. Usługodawca zachowuje prawa do własnych metod, narzędzi, szablonów i oprogramowania.
\end{enumerate}

\section*{Część IV. Tryb Fully private (offline)}

\subsection*{§13 Istota trybu}
\begin{enumerate}
\item W trybie Fully private Materiały są przetwarzane wyłącznie na komputerze bez połączenia z internetem (modele działają lokalnie), a dane nie trafiają do żadnej usługi chmurowej ani do zewnętrznego dostawcy AI.
\item Tryb realizuje się w jednym z wariantów wskazanych w ofercie: (A) na miejscu u Klienta, na urządzeniu Private AI Node lub urządzeniu Usługodawcy przywiezionym do Klienta, w obecności jego przedstawiciela; (B) w siedzibie Usługodawcy na wydzielonym komputerze offline, do którego Materiały dostarcza się na szyfrowanym nośniku.
\item Przetwarzanie odbywa się w [Polsce / miejscu wskazanym w ofercie] w obrębie Europejskiego Obszaru Gospodarczego.
\end{enumerate}

\subsection*{§14 Przekazanie, zabezpieczenia i usunięcie}
\begin{enumerate}
\item W wariancie B Materiały przekazuje się na szyfrowanym nośniku (hasło podaje Klient osobnym kanałem, np. telefonicznie) osobiście lub kurierem z potwierdzeniem odbioru. Klient odpowiada za bezpieczeństwo nośnika do chwili wydania i po jego zwrocie.
\item Komputer offline ma szyfrowany dysk, wyłączone interfejsy sieciowe i dostęp wyłącznie dla Usługodawcy. Usługodawca nie korzysta z podwykonawców przy przetwarzaniu Materiałów w tym trybie.
\item Wyniki przekazuje się na nośniku, w postaci wydruku lub na urządzeniu Klienta (wariant A). Kanał przekazania raportu zapisuje się w ofercie.
\item Po wydaniu raportu i potwierdzeniu odbioru przez Klienta, najpóźniej w ciągu [7] dni, Usługodawca trwale usuwa Materiały, kopie robocze i raport z komputera (bezpieczne wymazanie lub kasowanie kryptograficzne) i potwierdza to Klientowi pisemnie. W wariancie A usunięcie odbywa się w obecności przedstawiciela Klienta. Na pisemne życzenie Klienta Usługodawca przechowuje Materiały dłużej, nie dłużej niż [30] dni, na czas ewentualnych poprawek.
\item Usługodawca nie tworzy kopii zapasowych Materiałów poza szyfrowanymi kopiami roboczymi na komputerze offline, które usuwa na zasadach z ust. 4.
\item Naruszenie ochrony danych Usługodawca zgłasza Klientowi bez zbędnej zwłoki, nie później niż w ciągu 24 godzin od jego stwierdzenia.
\end{enumerate}

\subsection*{§15 Ograniczenia trybu}
\begin{enumerate}
\item Tryb Fully private wykorzystuje modele otwarte uruchamiane na sprzęcie lokalnym; ich jakość i szybkość mogą być inne niż w największych modelach chmurowych. Czas realizacji jest zwykle dłuższy niż w trybie Cloud, a cena wyższa, co wynika z oferty.
\item Brak zdalnego dostępu Usługodawcy do danych Klienta. Poprawki i dodatkowe pytania do raportu wymagają ponownego przekazania Materiałów, jeśli zostały już usunięte.
\end{enumerate}

\section*{Część V. Tryb Cloud (online)}

\subsection*{§16 Istota trybu}
\begin{enumerate}
\item W trybie Cloud Materiały są przetwarzane z użyciem zewnętrznych usług AI i infrastruktury chmurowej (dostawców), dzięki czemu analiza jest zwykle szybsza, a modele mocniejsze niż w trybie Fully private.
\item Materiały opuszczają więc komputer Usługodawcy i trafiają do dostawców wskazanych w Załączniku nr 1 (lista podprocesorów). Klient wybiera ten tryb świadomie.
\end{enumerate}

\subsection*{§17 Dostawcy i podpowierzenie}
\begin{enumerate}
\item Usługodawca korzysta tylko z dostawców, z którymi ma zawartą umowę powierzenia (art. 28 RODO) i którzy zobowiązali się, że dane Klienta nie są używane do trenowania modeli, a ich przechowywanie jest ograniczone do niezbędnego minimum (w miarę możliwości: brak przechowywania lub krótka retencja).
\item Klient wyraża ogólną zgodę na korzystanie z dostawców z Załącznika nr 1. O zamiarze dodania lub zmiany dostawcy Usługodawca informuje Klienta z wyprzedzeniem min. [14] dni; Klient może sprzeciwić się pisemnie, a wtedy zlecenie realizuje się w trybie Fully private lub umowa wygasa bez kar.
\item Jeśli dostawca przetwarza dane poza Europejskim Obszarem Gospodarczym, Usługodawca zapewnia mechanizm z rozdziału V RODO (np. decyzja o adekwatności lub standardowe klauzule umowne) i informuje o nim w Załączniku nr 1. Gdy to możliwe, wybiera się regiony w EOG.
\item Usługodawca odpowiada wobec Klienta za działania dostawców jak za własne, w granicach § 22.
\end{enumerate}

\subsection*{§18 Przekazanie, retencja i usunięcie}
\begin{enumerate}
\item Materiały przekazuje się wyłącznie bezpiecznym kanałem uzgodnionym w ofercie (np. szyfrowany udział plików z dostępem na zaproszenie i wygaśnięciem linku, szyfrowany plik z hasłem podanym osobno). Nie wolno używać zwykłego e-maila.
\item Dane o zdrowiu, inne dane szczególnych kategorii oraz materiały objęte tajemnicą zawodową mogą być przetwarzane w trybie Cloud tylko, jeśli Klient wyraźnie potwierdzi to w ofercie po zapoznaniu się z listą dostawców. W przeciwnym razie Klient je anonimizuje albo wybiera tryb Fully private.
\item Usługodawca usuwa Materiały, kopie robocze i dane z systemów własnych w ciągu [14] dni od przekazania raportu i potwierdza to pisemnie, a u dostawców zgodnie z ich okresem retencji wskazanym w Załączniku nr 1. Na pisemne życzenie Klienta usunięcie następuje wcześniej.
\item Naruszenie ochrony danych Usługodawca zgłasza Klientowi bez zbędnej zwłoki, nie później niż w ciągu 24 godzin od stwierdzenia.
\end{enumerate}

\subsection*{§19 Ograniczenia trybu}
\begin{enumerate}
\item Mimo zabezpieczeń tryb Cloud niesie wyższe ryzyko niż tryb offline (dostęp dostawców, ich incydenty, zmiany polityk). Klient przyjmuje to ryzyko, wybierając ten tryb.
\item Dostępność usług dostawców nie leży w kontroli Usługodawcy; przerwy po stronie dostawcy mogą przesunąć termin, bez odpowiedzialności Usługodawcy, o ile dochował należytej staranności.
\end{enumerate}

\section*{Część VI. Płatności, reklamacje, odpowiedzialność i postanowienia końcowe}

\subsection*{§20 Wynagrodzenie i płatności}
\begin{enumerate}
\item Ceny wynikają z oferty i są podawane netto; do cen doliczany jest podatek VAT według obowiązującej stawki. Usługodawca wystawia faktury, w tym w KSeF.
\item Jeśli oferta nie stanowi inaczej, Klient wpłaca zaliczkę [50]\% przed rozpoczęciem pracy, a resztę w ciągu [14] dni od faktury wystawionej po przekazaniu raportu.
\item Rozmowa wstępna i wycena są bezpłatne.
\item Prace spoza zakresu oferty (np. dodatkowe Materiały, zmiana zakresu, ponowna analiza po usunięciu danych) wymagają nowej wyceny i akceptacji Klienta.
\end{enumerate}

\subsection*{§21 Reklamacje}
\begin{enumerate}
\item Klient zgłasza zastrzeżenia do raportu pisemnie na [e-mail] w ciągu [14] dni od jego otrzymania, opisując niezgodność z ofertą.
\item Usługodawca odpowiada w ciągu [14] dni. Uzasadniona reklamacja skutkuje bezpłatną poprawą raportu w uzgodnionym terminie; jeśli poprawa nie jest możliwa, Usługodawca obniża wynagrodzenie proporcjonalnie.
\item Reklamacją nie jest niezgoda co do oceny merytorycznej lub prawnej sprawy ani błąd modelu AI, który Klient mógł wykryć przy weryfikacji z źródłami zgodnie z § 8 ust. 3, chyba że raport nie spełnia zakresu z oferty.
\end{enumerate}

\subsection*{§22 Odpowiedzialność}
\begin{enumerate}
\item Odpowiedzialność Usługodawcy z tytułu danego zlecenia, z wyłączeniem szkody wyrządzonej umyślnie, ogranicza się do wysokości wynagrodzenia zapłaconego za to zlecenie. Usługodawca nie odpowiada za utracone korzyści ani szkody pośrednie.
\item Usługodawca nie odpowiada za szkody wynikające z braku weryfikacji raportu przez Klienta, z przekazania danych niezgodnie z Regulaminem, z niedokładnych lub niekompletnych Materiałów oraz z decyzji podjętych przez Klienta na podstawie raportu.
\item Ograniczenia z ust. 1–2 nie naruszają odpowiedzialności Usługodawcy wobec Klienta za naruszenie obowiązków poufności i ochrony danych w zakresie, w jakim prawo nie pozwala jej wyłączyć ani ograniczyć. [Do decyzji z prawnikiem: odrębny, wyższy limit dla naruszenia poufności i danych.]
\item Strony nie odpowiadają za niewykonanie zobowiązania spowodowane siłą wyższą.
\end{enumerate}

\subsection*{§23 Rezygnacja i rozwiązanie}
\begin{enumerate}
\item Klient może zrezygnować ze zlecenia przed rozpoczęciem pracy; zaliczka jest zwracana w ciągu 14 dni, pomniejszona o udokumentowane koszty poniesione do tego czasu.
\item Po rozpoczęciu pracy Klient płaci za pracę wykonaną do dnia rezygnacji, proporcjonalnie do postępu.
\item Usługodawca może odstąpić od zlecenia z ważnej przyczyny, w szczególności gdy Materiały zostały przekazane niezgodnie z Regulaminem lub ich przetwarzanie byłoby niezgodne z prawem. Wtedy Usługodawca zwraca zaliczkę za pracę niewykonaną i usuwa Materiały zgodnie z właściwym trybem.
\end{enumerate}

\subsection*{§24 Postanowienia końcowe}
\begin{enumerate}
\item Informacje o przetwarzaniu danych osobowych przez Usługodawcę (kontakt, rozmowy, nagrania, rozliczenia) zawiera Załącznik nr 1.
\item Usługodawca może zmienić Regulamin; zmiana nie dotyczy zleceń zaakceptowanych przed jej ogłoszeniem. O zmianie informuje się na stronie i e-mailem aktywnych Klientów z wyprzedzeniem [14] dni.
\item Do umów stosuje się prawo polskie. Spory rozstrzyga sąd właściwy dla siedziby Usługodawcy.
\item Regulamin obowiązuje od dnia [data].
\end{enumerate}

\clearpage
\section*{Załącznik nr 1: klauzula informacyjna (art. 13 RODO) i odbiorcy danych}

\begin{tabularx}{\linewidth}{@{}>{\bfseries}p{0.27\linewidth}X@{}}
\toprule
Pytanie & Odpowiedź \\
\midrule
Administrator & Konrad Cinkusz, Inference Polska, [adres], e-mail [adres] \\
Cele i podstawy & kontakt i wycena (art. 6 ust. 1 lit. b i f); nagrywanie rozmów bezpłatnych, notatka i analiza AI (zgoda, art. 6 ust. 1 lit. a); realizacja umów i rozliczenia (art. 6 ust. 1 lit. b i c); dochodzenie lub obrona roszczeń (art. 6 ust. 1 lit. f) \\
Kategorie danych & imię i nazwisko, dane kontaktowe, nazwa firmy, wizerunek i głos w nagraniu, treść rozmowy \\
Okres przechowywania & nagranie do [90] dni; notatka do [12] miesięcy od ostatniego kontaktu; dokumenty księgowe zgodnie z przepisami (5 lat od końca roku podatkowego); dane do roszczeń do ich przedawnienia \\
Odbiorcy & Google (Google Meet, Poczta) jako dostawca spotkań i poczty; [narzędzie transkrypcji / analizy AI lub: wyłącznie lokalnie]; biuro rachunkowe [nazwa]; dostawca KSeF/faktur [nazwa] \\
Transfer poza EOG & [dotyczy Google i ewentualnie dostawców AI: mechanizm, np. EU-U.S. Data Privacy Framework lub standardowe klauzule umowne] \\
Prawa & dostęp, sprostowanie, usunięcie, ograniczenie, sprzeciw, przeniesienie, cofnięcie zgody w każdej chwili, skarga do Prezesa UODO \\
Zautomatyzowane decyzje & brak \\
Dobrowolność & podanie danych jest dobrowolne, ale potrzebne do kontaktu i umowy; zgoda na nagrywanie jest dobrowolna \\
\bottomrule
\end{tabularx}

\subsection*{Dostawcy w trybie Cloud (§ 17)}
\begin{tabularx}{\linewidth}{@{}XXXXX@{}}
\toprule
Dostawca & Cel & Region przetwarzania & Mechanizm transferu & Retencja u dostawcy \\
\midrule
{}[nazwa] & [model AI / hosting] & [EOG / USA] & [mechanizm] & [np. 0 dni / 30 dni] \\
{}[nazwa] & [przechowywanie plików] & [region] & [mechanizm] & [okres] \\
\bottomrule
\end{tabularx}

\section*{Załącznik nr 2: treść informacji o nagrywaniu (do zaproszenia Meet i do odczytania na początku)}

\begin{itemize}
\item Tekst do zaproszenia i e-maila: „Rozmowa odbędzie się w Google Meet i, za Twoją zgodą, zostanie nagrana. Nagranie służy mi do sporządzenia notatki i analizy z pomocą AI; usuwam je po [90] dniach, nie używam do trenowania modeli i nie udostępniam osobom trzecim poza opisanymi w klauzuli. Możesz odmówić nagrywania bez wpływu na rozmowę. Proszę nie podawać danych pacjentów ani klientów. Pełna klauzula: [link].”
\item Formuła ustna na starcie: „Zanim zaczniemy: nagrywam tę rozmowę do własnej notatki i analizy z AI. Czy wyrażasz zgodę? Możesz odmówić, a ja będę wtedy tylko notował.” Potem włącz nagrywanie dopiero po usłyszeniu zgody od wszystkich uczestników.
\end{itemize}

\end{document}
